\section{Bj\"ork's approach to existence of b-functions}

\begin{theorem}[Bj\"ork]\label{Theorem 8.4}
    Let $M$ be a holonomic $\mathcal{D}_X$-module with a good filtration $\Omega_{\bullet} M$ over $V_{\bullet} \mathcal{D}_X$. There exists a non zero $b(s) \in \mathbb{C}[s]$ such that
    \[
        b(t\partial_t + k) \text{ kills } \frac{\Omega_k M}{\Omega_{k-1} M}, \quad \forall k \in \mathbb{Z}.
    \]
    The monic minimal degree $b(s)$ is called the $\textbf{$b$-function}$ of $M$ along $Y$.
\end{theorem}
\begin{remark}
    In the handwritten notes provided by the lecturer, says $b(t\partial_t \textcolor{red}{ - } k) \text{ kills } \frac{\Omega_k M}{\Omega_{k-1} M}$,
    but this is not compatible with the increasing filtration setting in Rmk.~\ref{rmk:increasing_V_filtration}.
\end{remark}


\begin{lemma}\label{lemma 8.5}
    Let $Y$ be a smooth variety over $\mathbb{C}$, $N$ a holonomic $\mathcal{D}_Y$-module, and $\theta\in \mathrm{End}_{\mathcal{D}_Y}(N)$. Then there exists $b(s)\in\mathbb{C}[s]$ such that $b(\theta)=0$.
\end{lemma}
\begin{proof}
    Since $N$ is holonomic, we obtain that
    \[
        \mathrm{Ch}(N) = \bigcup_{i} T^*_{Y_i} Y,
    \]
    with $Y_{i}$ locally closed.

    Take a $Y_i$ of maximal dimension and take an (affine) open subset $U \subset Y$
    such that $Y_{i}^{\circ} \hookrightarrow U$ is a closed embedding
    with $Y_{i}^{\circ}$ smooth, $\overline{Y_i^{\circ}} = Y_{i}$, $U \cap Y_i = Y_i^\circ$
    and $U \cap (\bigcup_{j\neq i}Y_j) = \varnothing$. Then
    \[
        \mathrm{Ch}(N|_U) = T^*_{Y_i^\circ} U.
    \]

    Since $Y^\circ$ is smooth, locally we can split algebraic coordinates to use Kashiwara's equivalence \ref{thm:kashiwara_equiv}, by shrinking $U$:
    \[
        N|_U = i_{+}M = M[\partial_1, \ldots, \partial_k].
    \]

    Consider $U^{\mathrm{an}} \to U$, the analytification functor, since $Y^\circ \hookrightarrow U$ is a closed embedding,
    by \cite[Proposition~4.7.2.(ii)]{hotta2008d} we have $i_{+}(M)^{\mathrm{an}} \to i^{\mathrm{an}}_{+}(M^{\mathrm{an}})$ is an isomorphism.

    We use Riemann-Hilbert correspondence on $(N|_U)^{\mathrm{an}}$:
    \begin{align*}
        (N|_U)^{\mathrm{an}} & \cong (i_{+}M)^{\mathrm{an}}                                                                                                       &  & \text{(By \ref{thm:kashiwara_equiv})}                                                                                   \\
                             & \cong i^{\mathrm{an}}_{+}(M^{\mathrm{an}})                                                                                                                                                                                                                      \\
                             & \cong M^{\mathrm{an}}[\partial_1, \ldots, \partial_k]                                                                              &  & \text{\begin{tabular}[t]{@{}l@{}}($i_+$ of a closed embedding behaves\\ the same on variety and manifold)\end{tabular}} \\
                             & \cong \mathcal{O}^{\mathrm{an}}_{Y_i^\circ} \otimes_{\mathbb{C}} L \otimes_{\mathbb{C}} \mathbb{C}[\partial_1, \ldots, \partial_k] &  & \text{(By \ref{cor:kashiwara_equiv_conormal}, $M$ is an integrable connection, use \ref{thm:basic-RH-corr})}
    \end{align*}
    Then we get
    \[
        (N|_U)^{\mathrm{an}} = \mathcal{O}^{\mathrm{an}}_{Y_i^\circ} \otimes L[\partial_1, \ldots, \partial_k]
    \]
    with $L$ a local system on $Y_i^\circ$.

    For the name \emph{Riemann-Hilbert correspondence}, one may think it is a \textbf{categorical equivalence}, that means there is not only a correspondence of objects but also a correspondence of \textbf{morphisms}.

    That is \cite[Theorem~4.2.4.]{hotta2008d}'s $H^{-\mathrm{d}_X}(\mathrm{DR}(-))$ ($X$ a complex manifold),
    which means for a morphism between two holomorphic $\mathcal{D}_X$-modules $f: M_1 \to M_2$ locally free over $\mathcal{O}_X$ (holomorphic),
    with respective flat sections
    \[
        L_1 := \Ker(\nabla_{M_1}) \quad \text{and}\quad L_2 := \Ker(\nabla_{M_2}) \quad \text{(see \ref{rmk:basic-RH-corr})},
    \]
    we have a morphism between their corresponding local systems
    \[
        H^{-\mathrm{d}_X}(\mathrm{DR}(f)): L_1 \to L_2
    \]
    which is just a restriction of $f$.

    Recall \ref{prop:classification-of-local-system} says if $X$ is connected, we have an another categorical equivalence
    to identify local system and finite diemsional representation of $\pi_1(X,x)$ (take an arbitrary base point $x$), morphisms of representations
    is called \textbf{intertwining operator} with respect to the action of $\pi_1(X,x)$.

    They compose a third categorical equivalence, since categorical equivalence is fully faithful, we have
    \begin{align*}
        \mathrm{End}_{\mathcal{D}_U}(N|_U) & = \mathrm{End}_{\mathcal{D}_{Y^\circ}}(M)                                     &  & \text{(Kashiwara equivalence is also a categorical equivalence)}          \\
                                           & \subseteq \mathrm{End}_{\mathcal{D}^{\mathrm{an}}_{Y^\circ}}(M^{\mathrm{an}})                                                                                \\
                                           & = \mathrm{End}_{\mathrm{LocSys}_{X}}(L)                                       &  & \text{(By \cite[Theorem~4.2.4.]{hotta2008d})}                             \\
                                           & = \mathrm{End}_{\pi_1(X,x)}(L_x)                                              &  & \text{(By \ref{prop:classification-of-local-system}, and $L_x$ is fiber)} \\
                                           & \subseteq \mathrm{End}_{\C}(L_x)                                              &  & \text{(Definition of intertwining operator)}
    \end{align*}
    Therefore, by Cayley-Hamilton theorem, there exists $b_i(s) \in \mathbb{C}[s]$ such that $b_i(\theta|_{U}) = 0$.

    Before we do induction, recall a fact that for a smooth variety $Y$, a irreducible closed subset $X$, we have
    \[
        T^*_X Y = \overline{T^*_{X\cap U} U} \subseteq T^*Y \quad \text{(for affine open $U \subset Y$ such that $X\cap U \subseteq X^{\mathrm{sm}}$ dense in $X$)}
    \]
    is also irreducible, by the fact that
    \[
        T^*_{X \cap U} U \subseteq T^*U.
    \]
    is closed irreducible in $T^*U$.

    Now, take  $N_1 = b_1(\theta)N$.
    Since $b_1(\theta|_U) = 0$, we have $N_1|_U = 0$, which implies $\mathrm{Ch}(N_1) \cap T^*U = \varnothing$.

    By Bernstein's inequality \ref{thm:Bernstein-inequality},
    $N_1$, as a submodule of a holomonic $\D_Y$-module is also a holomonic $\D_Y$-module.
    So irreducible component of $\mathrm{Ch}(N_1)$ has dimension $n = \dim Y$,
    and is contained in some irreducible component of $\mathrm{Ch}(N) = \bigcup_k T^*_{Y_k} Y$.

    Since all irreducible component of a holomonic $\D_Y$-module has dimension $n = \dim Y$ (see Cor.~\ref{cor:holonomic-char-cycle}),
    and we have a fact that a proper closed subset of an irreducible variety has strictly smaller dimension,
    it follows that each irreducible component of $\mathrm{Ch}(N_1)$ must be equal to some $T^*_{Y_k} Y$.

    But it can't be $T^*_{Y_i} Y$, because $\mathrm{Ch}(N_1) \cap T^*U = \varnothing$.
    So we strictly reduce the number of components in the characteristic variety,
    so we are done by induction.


\end{proof}




\begin{theorem}\label{Theorem 8.6}
    Let $M$ be a finitely generated $\mathcal{D}_{X}$-module, and $\Omega_{\bullet}M$ a good filtration over $V_{\bullet}\mathcal{D}_{X}$. Then $j(\mathrm{gr}_{\bullet}^\Omega M)\ge j(M)$.
\end{theorem}
\begin{remark}
    In Thm.~\ref{theorem 7.12} we also proved $j(\gr_\bullet M)\ge j(M)$, but this time we have a different filtration, the $V$-filtration.

    We need the fact that $\gr_\bullet^{V}\D_X$ is Auslander regular, see Cor.~\ref{cor:Rees-ring-of-V-filtration-D-module}. So it can't replace the first proof in Thm.~\ref{theorem 7.12}.
\end{remark}
\begin{proof}
    If $\mathrm{gr}^\Omega_{\bullet}M=0$, there is nothing to prove, so we assume $\mathrm{gr}^\Omega_{\bullet}M\neq 0$ and $\ell=j(\mathrm{gr}_{\bullet}^{\Omega}M)<j(M)$.

    Then the $E_1$-page of the spectral sequence of $\check{F}^\bullet$
    introduced in Section~\ref{sect:more_on_filtered_rings},
    that converges to $\mathrm{Ext}_{\mathcal{D}_{X}}^{j}(M,\mathcal{D}_{X})$,
    satisfies:
    \[
        0 \rightarrow \operatorname{Ext}^{\ell}(\operatorname{gr}_{\bullet}^{\Omega} M, \operatorname{gr}^{V}_{\bullet}\mathcal{D}_X)
        \xrightarrow{\mathrm{d}_1}
        \operatorname{Ext}^{\ell+1}(\operatorname{gr}_{\bullet}^{\Omega} M, \operatorname{gr}^{V}_{\bullet}\mathcal{D}_X).
    \]
    Where we identify $\gr_\bullet \Hom_{\D_X}(F^\bullet, \D_X)$ with $\Hom_{\gr_\bullet^{V} \D_X}(\gr_\bullet F^\bullet, \gr_\bullet^{V} \D_X)$
    as in the proof of Cor.~\ref{corollary 7.11}.

    Since $\ell = j(\gr_\bullet^\Omega M)$, so we obtain that $E^{\ell}_{\bullet}(2)=\ker \mathrm{d}_1\subseteq\operatorname{Ext}^{\ell}(\operatorname{gr}_{\bullet}^{\Omega} M, \operatorname{gr}^{V}_{\bullet}\mathcal{D}_X)$ since no images map to $E^{\ell}_{\bullet}(1)$.

    Then we have the following exact sequence with the quotient $N=E^{\ell}_{\bullet}(1) /\ker \mathrm{d}_{1}\simeq \mathrm{d}_{1}(E_{\bullet}^{\ell}(1))\subseteq E^{\ell+1}_{\bullet}(1)$:
    \[
        0 \rightarrow E_{\bullet}^{\ell}(2) \rightarrow \operatorname{Ext}^{\ell}(\operatorname{gr}_{\bullet}^{\Omega} M, \operatorname{gr}^{V}_{\bullet}\mathcal{D}_X)
        \rightarrow N \rightarrow 0.
    \]
    In general, by induction, $\forall k\ge 1$, we also have such exact sequence:
    \[
        0 \rightarrow E_{\bullet}^{\ell}(k+1) \rightarrow E_{\bullet}^{\ell}(k) \rightarrow N_k \rightarrow 0.
    \]

    Since $j(M)>\ell$, and $\check{F}^\bullet$ is convergent (see Definition~\ref{def:convergent_spectral_sequence}),
    we have $E_{\bullet}^{\ell}(r) = \operatorname{gr}_{\bullet} \operatorname{Ext}^{\ell}(M, \mathcal{D}_X) = 0$ for some $r\gg 0$,
    and $E^{\ell}_{\bullet}(r)=\Ker \mathrm{d}_{r}$.

    By triangle inequality of grade number (one can get this inequality by a long exact sequence corresponds to the short exact sequence), we have
    \[
        j(E_{\bullet}^{\ell}(k)) \ge \min(j(E_{\bullet}^{\ell}(k+1)), j(N_k)).
    \]
    Note that $N_k$ is a subquotient of $E_{\bullet}^{\ell+1}(1) = \operatorname{Ext}^{\ell+1}(\operatorname{gr}_{\bullet}^{\Omega} M, \operatorname{gr}^{V}_{\bullet}\mathcal{D}_X)$.
    Since $\gr_\bullet^{V}\D_X$ is Auslander regular, we have $j(N_k) \ge \ell+1$.

    Since $E_{\bullet}^{\ell}(r) = 0$ for $r \gg 0$, we have $j(E_{\bullet}^{\ell}(r)) = \infty$.
    By induction (downwards on $k$), we obtain
    \[
        j(E_{\bullet}^{\ell}(1)) \ge \ell+1.
    \]
    However, $E_{\bullet}^{\ell}(1) = \operatorname{Ext}^{\ell}(\operatorname{gr}_{\bullet}^{\Omega} M, \operatorname{gr}^{V}_{\bullet}\mathcal{D}_X)$.
    Since $\ell = j(\operatorname{gr}_{\bullet}^{\Omega} M)$, we have $j(E_{\bullet}^{\ell}(1)) = \ell$ by Theorem~\ref{theorem 7.2}'s (2) and (4).
    This contradicts $j(E_{\bullet}^{\ell}(1)) \ge \ell+1$.
\end{proof}

\begin{proof}[proof of (\ref{Theorem 8.4})]
    \quad

    By Theorem~\ref{Theorem 8.6},
    $j(\mathrm{gr}_{\bullet}^\Omega M)\ge j(M)=n$.
    If $\mathrm{gr}_{\bullet}^\Omega M=0$, take $b(s)=1$.

    Since by Cor.~\ref{cor:Rees-ring-of-V-filtration-D-module} and filtration is good, we view $\mathrm{gr}_{\bullet}^\Omega M$ as a finitely generated$\gr_\bullet^{V}\D_X = \D_{T_Y X}$-module.
    And in the proof of Cor.~\ref{cor:Rees-ring-of-V-filtration-D-module}, we verified that $T_{Y}X$ is a smooth variety of dimension $n$ over $\mathbb{C}$.
    So if $\gr_\bullet^{\Omega}M \neq 0$, then $\mathrm{dim}(\mathrm{Ch}(\mathrm{gr}^\Omega_{\bullet}M))\ge n$ by Thm.~\ref{thm:Bernstein-inequality}.

    Since $\mathrm{gr}_{\bullet}^V\mathcal{D}_{X}$ is a Auslander regular ring,
    by Thm.~\ref{theorem 7.12},
    we obtain that $j(\mathrm{gr}_{\bullet}^\Omega M)+\mathrm{dim}(\mathrm{Ch}(\mathrm{gr}_{\bullet}^\Omega M))=2n=2\,\mathrm{dim}T_{Y}X$, then $\mathrm{gr}^\Omega_{\bullet} M$ is holonomic over $T_Y X$.

    Now, take
    \[\theta=\bigoplus \limits_{i}(t\partial_{t}+i) : \gr_\bullet^\Omega M \to \gr_\bullet^\Omega M\quad \text{ for } t\partial_t + i: \gr_i^\Omega M \to \gr_i^\Omega M\]
    since one can verify that for $P\in\D_X$
    \[P \in V_{i}\D_X \Leftrightarrow P = p(x,\partial_x) t^{\alpha}\partial_t^\beta\quad \text{ for } p(-,-) \text{ a polynomial, and }\beta - \alpha \le i\]
    which implies
    \[[t\partial_t, P] = -i P \text{ for } P\in\gr_i^{V}\mathcal{D}_X,\]
    so we have $\theta\in\mathrm{End}_{\mathcal{D}_{T_{Y}X}}(\mathrm{gr}_{\bullet}^\Omega M)$.
    Indeed, for $u \in \gr_j^\Omega M$ and $P \in \gr_k^{V} \D_X$, we calculate
    \[
        \theta(Pu) = (t\partial_t + j + k) Pu = (P t\partial_t - k P + j + k) u = P(t\partial_t + j) u = P \theta(u).
    \]

    By Lemma \ref{lemma 8.5}, there exists $b(s)\in\mathbb{C}[s]$ s.t. $b(\theta)=0$,
    In particular, $b(t\partial_{t}+k)$ kills $\frac{\Omega_{k}M}{\Omega_{k-1}M}$ for all $k$.
\end{proof}

\begin{remark}\label{rmk:b_function_of_C[x,1/x]e^1/x}
    It is possible that $\mathrm{gr}_{\bullet}^\Omega M=0$. In fact, take
    $M := \mathcal{D}_{\Spec{\C[t]}} \big/ \mathcal{D}_{\Spec{\C[t]}}\langle t^2\partial_t + 1 \rangle = \mathcal{D}_{X}\cdot e^{\frac{1}{t}}$,
    for $e^{\frac{1}{t}}$ represents the canonical section $[1]\in\mathcal{D}_{\Spec{\C[t]}} \big/ \langle t^2\partial_t + 1 \rangle$.

    We note that $t$ also acts bijectively on $M$ (just as Example~\ref{ex:first-examples-of-d-modules}),
    so $M=M[t^{-1}]  = \mathbb{C}[t, \frac{1}{t}]e^{\frac{1}{t}}$,
    take $\Omega_{k}M=V_{k}\mathcal{D}_{X}\cdot e^{\frac{1}{t}}$,
    which is a good filtration.
    Since $t\partial_{t}e^{\frac{1}{t}}= -\frac{1}{t} e^{\frac{1}{t}}$, we have $\Omega_{k}M=M$ and $b(s)=1$.
\end{remark}

\begin{theorem}[Kashiwara-Malgrange]\label{Theorem 8.7}
    Let $M$ be a holonomic $\mathcal{D}_{X}$-module, there exists unique good filtration $V_{\bullet}M$ and $b(s)\in\mathbb{C}[s]$ s.t.
    \begin{enumerate}
        \item $\mathrm{Re}(\text{roots of }b(s))\subseteq[0,1)$;
        \item $b(t\partial_{t}+k)$ kills $\frac{V_{k}M}{V_{k-1}M}$.
    \end{enumerate}
    And the filtration is called $V$-filtration.
\end{theorem}
\begin{proof}
    This is Kashiwara's trick:
    First, set $\Gamma_{n}=\Omega_{n+k}$. Note that $b(t\partial_{t}+n+k)$ kills $\Gamma_{n} /\Gamma_{n-1}$ for all $n$,
    which means the roots of the $b$-function are shifted by $k$.
    Thus by taking $k\gg 0$, after replacing $b(s)$ by $s\mapsto b(s+k)$, we can assume $\mathrm{Re}(\lambda)\le 0$ for all roots $\lambda$ of $b(s)$.

    Next, suppose that $b(s)$ has a root $\alpha$ whose real part is $< 0$.
    Set $\Gamma_{n}'=\Gamma_{n-1}+(t\partial_{t}+n-\alpha)\Gamma_{n}$.
    If $\beta(s) = b(s)/(s-\alpha)$, then we verify that $\beta(s+n)(s+n-\alpha-1)$ annihilates $\Gamma'_{n} / \Gamma'_{n-1}$ for all $n$.

    Denote $s = t\partial_t$. Let $u' \in \Gamma'_{n}$. We can write $u' = v + (s+n-\alpha)w$ with $v \in \Gamma_{n-1}$ and $w \in \Gamma_{n}$.
    We calculate the action of $Q(s) = \beta(s+n)(s+n-\alpha-1)$:
    \[
        Q(s)u' = \beta(s+n)(s+n-\alpha-1)v + \beta(s+n)(s+n-\alpha-1)(s+n-\alpha)w.
    \]
    For the first term involving $v \in \Gamma_{n-1}$: note that $(s+n-\alpha-1)v \in (s+(n-1)-\alpha)\Gamma_{n-1} \subseteq \Gamma'_{n-1}$.
    For the second term involving $w \in \Gamma_{n}$: we use the fact that polynomials in $s$ commute, so
    \[
        \beta(s+n)(s+n-\alpha-1)(s+n-\alpha)w = (s+n-\alpha-1) \beta(s+n)(s+n-\alpha) w = (s+n-\alpha-1) b(s+n) w.
    \]
    By the hypothesis on $\Gamma_\bullet$, we have $b(s+n)w \in \Gamma_{n-1}$. Let $z = b(s+n)w$.
    Then $(s+n-\alpha-1)z \in (s+(n-1)-\alpha)\Gamma_{n-1} \subseteq \Gamma'_{n-1}$.
    Thus, $Q(s)u' \in \Gamma'_{n-1}$.

    Continue doing this if $\alpha$ has multiplicity, we will finally move a root $\alpha$ of $b(s)$ to $\alpha+1$.
\end{proof}

\begin{definition}
    $M$ is a holomonic module, then:

    \begin{enumerate}
        \item  $M$ is called $Q\textbf{-specializable}$ along $Y$ if the $b$-function has all its roots in $Q$, where $Q \subseteq \mathbb{C}$ is a subfield.

              Since any subfield of $\C$ must contain its \textbf{prime subfield} $\mathbb{Q}$, and $b(t\partial_t + k)$ kills $V_kM/V_{k-1}M$,
              the eigenvalues $\lambda$ of $t\partial_t$ on $V_kM/V_{k-1}M$ kills its eigenvectors by multiplying $b(\lambda + k)$.

              So $\lambda$ differs from the roots of $b(s)$ by integers.
              Thus, if all roots of $b(s)$ are in $Q$, then all eigenvalues of $t\partial_t$ are in $Q$ (as $\mathbb{Z} \subseteq \mathbb{Q} \subseteq Q$).

        \item If $M$ is $\mathbb{Q}$-specializable,
              then for all $\beta \in \mathbb{Q}$,
              take $k = \lceil \beta \rceil$,
              then define
              $$V_{\beta} M = \pi^{-1} \left( \bigoplus_{\alpha \leq \beta} \left( \frac{V_k M}{V_{k-1} M} \right)^{\alpha} \right),$$
              and
              $$V_{<\beta} M = \pi^{-1} \left( \bigoplus_{\alpha < \beta} \left( \frac{V_k M}{V_{k-1} M} \right)^{\alpha} \right),$$
              where $\left( \frac{V_k M}{V_{k-1} M} \right)^{\alpha}$ denotes the $\alpha$-generalized eigenspace of $\frac{V_k M}{V_{k-1} M}$
              with respect to the $t\partial_t$-action, and $\pi: V_k M \rightarrow \frac{V_k M}{V_{k-1} M}$;
              Then we get a $\mathbb{Q}$-indexed $V$-filtration s.t.
              $$\mathrm{gr}_{\beta}^{V} M = \frac{V^{\beta} M}{V^{<\beta} M} = \left( \frac{V_k M}{V_{k-1} M} \right)^{\beta}$$
              with $t\partial_t - \beta$ acts on $\mathrm{gr}_{\beta}^{V} M$ nilpotently.
    \end{enumerate}
\end{definition}
\begin{example}
    Let $X=\mathbb{A}^{1}_{\mathbb{C}}$ and $M = \mathbb{C}\left[ t, \frac{1}{t} \right]e^{\frac{1}{t}}$ an irregular holonomic $\mathcal{D}$-mod.

    Consider $V$-filtration (of $\D_X$) along $t = 0$. The $V$-filtration on $M$ exists by Theorem \ref{Theorem 8.7}.

    Recall Rmk.~\ref{rmk:b_function_of_C[x,1/x]e^1/x}. The filtration $\Omega_\bullet M$ is defined by $\Omega_i M := V_i \mathcal{D}_{X} \cdot e^{\frac{1}{t}}$ for $i \in \mathbb{Z}$.
    We have $\Omega_i M = M$ for all $i\in \mathbb{Z}$.

    This filtration gives a $b$-function $b(s) = 1$. Which has all its roots in $[0,1)$, so $\Omega_\bullet M$ is exactly the $V$-filtration on $M$.
\end{example}

\begin{theorem}
    If M is a regular holonomic $\mathcal{D}_X$-mod, then its V-filtration is not trivial. In particular, its b-function is not $1$.
\end{theorem}

\begin{proof}[Sketch of proof]
    First, each $V_i M$ is a $V_0 \mathcal{D}_X$-module and $V_0 \mathcal{D}_X = \mathcal{D}_X^{\log}$ w.r.t. $(t = 0)$. Since $V_i M$ has no $t$-torsion for all $i < 0$, $\mathrm{Ch}(V_i M)$ has no component over $(t = 0)$, for all $i < 0$.

    Further, we obtain that $\mathrm{Ch}^{\log}\left(\frac{V_i M}{V_{i-1} M}\right) = \mathrm{Ch}^{\log}\left(\frac{V_i M}{t V_i M}\right)= \mathrm{Ch}^{\log}(V_i M)\big|_{(t=0)}$, thus $\frac{V_i M}{V_{i-1} M} \neq 0.$

\end{proof}

\begin{definition}[$\D$-module theoretic direct image of a graph embedding]
    Firstly we introduce the definition of graph embedding: 1

\end{definition}

\begin{example}\label{Example 8.11}

    Let $X = \mathbb{A}_{\mathbb{C}}^n$ and $Y = \mathbb{A}_{\mathbb{C}}^n \times \mathbb{A}_{\mathbb{C}}^1$, here $\mathbb{A}^{1}_{\mathbb{C}}$ is with coordinate $t$.
    For $f \in \mathbb{C}[x_1, \cdots, x_n]$, Consider $M = \mathbb{C}[x_1, \cdots, x_n]_f[s] f^s$.

    Define $t M = f M$ and then $t^{-1} M = f^{-1} M$.  Take $s = -t\partial_t$, then $\partial_t = -t^{-1} s$, the $\partial_{t}$-action is defined, and further, $M$ is a $\mathcal{D}_Y = \mathcal{D}_X \langle t, \partial_t \rangle$-module.

    Then $N := \mathbb{C}[x_1, \cdots, x_n][\partial_t] f^s = \mathbb{C}[x_1, \cdots, x_n][\partial_t] f^s= \mathbb{C}[x_1, \cdots, x_n][t^{-1} s] f^s$, and $\partial_{x_i} f^s = \partial_{x_i}(f) \cdot \frac{1}{f} s f^s= \partial_{x_i}(f) \cdot (-\partial_t) f^s$, thus $N$ is a $\mathcal{D}_Y$-submodule of $M$. Furthermore, we can check that $\mathrm{Ch}(N)$ = the conormal bundle of the graph of $X$, which indicates that $N$ is a holonomic $\mathcal{D}_{Y}$-module.

    We consider $V_{i}\mathcal{D}_{Y}\cdot f^{s}$ forming a good filtration of $N$: $$V_0 \mathcal{D}_Y \cdot f^s = \mathcal{D}_X \langle t, s \rangle f^s = \mathcal{D}_X[s] f^s$$Similarly, $V_{-1} \mathcal{D}_Y \cdot f^s = \mathcal{D}_X[s] f^{s+1}$.

    By Theorem \ref{Theorem 8.4}, there exists $b(s) \in \mathbb{C}[s]$ s.t.$$
        b(s) \cdot \frac{V_0 \mathcal{D}_Y \cdot f^s}{V_{-1} \mathcal{D}_Y \cdot f^s} = b(s) \cdot \frac{\mathcal{D}_X[s] f^s}{\mathcal{D}_X[s] f^{s+1}} = 0.$$which is to say there exists $P(s)\in \mathcal{D}_{X}[s]$ such that $b(s) f^s = P(s) \cdot f^{s+1}$

    In fact, $N = \iota_{f,+} (\mathbb{C}[x_1, \ldots, x_n])$, the $\mathcal{D}$-module pushforward, where $$\iota_f: X \to Y = X \times \mathbb{A}^{1}_{\mathbb{C}}\atop x \mapsto (x, f(x))$$the graph embedding. So $N$ is regular holonomic, and $b(s) \neq 1$.
\end{example}

\begin{remark}
    One can replace $\mathbb{C}[x_1, \ldots, x_n]$ by any regular holonomic $\mathcal{D}_X$-module, to get similar results in general.
\end{remark}
